\documentclass[a4paper,12pt]{article}
\usepackage[T2A]{fontenc}
\usepackage[utf8]{inputenc}
\usepackage{cmap}
\usepackage[english, russian]{babel}
\usepackage{wrapfig}
\usepackage[backend=biber,style=gost-numeric]{biblatex} % стиль ГОСТ
\addbibresource{literatura.bib}
\usepackage{misccorr} % в заголовках появляется точка, но при ссылке на них ее нет
\usepackage{amssymb,amsfonts,amsmath,amsthm}  
\usepackage{indentfirst}
\usepackage[usenames,dvipsnames]{color} 
\usepackage[unicode,hidelinks]{hyperref}
% \hypersetup{%
%     pdfborder = {0 0 0}
\usepackage{multirow}
\usepackage{makecell,multirow} 
\usepackage{ulem}
%\usepackage{hyperref}
%\usepackage{graphicx,wrapfig}
%\renewcommand\epsilon{\varepsilon}
% \renewcommand\epsilon{\varepsilon}
%\graphicspath{{img/}}
\usepackage{geometry}
\geometry{left=1cm,right=1cm,top=2cm,bottom=2cm,bindingoffset=0cm,headheight=15pt}
\usepackage{fancyhdr} 
\linespread{1.2} 
\frenchspacing 
\renewcommand{\labelenumii}{\theenumii)} 
% \usepackage{caption}
%%%%%%%%%%%%%%%%%%%%%%%%%%%%%%%%%%%%%%%%%%%%%%%%%%%%%%%%%%%%%%%%%%%%%%%%%%%%%%%
%%%%%%%%%%%%%%%%%%%%%%%%%%%%%%%%%%%%%%%%%%%%%%%%%%%%%%%%%%%%%%%%%%%%%%%%%%%%%%%

\def\labauthor{Сарафанов И.Г.}
\def\labauthors{\labauthor}
\def\labnumber{127}
\def\labtheme{Таблица производных и интегралов}

%%%%%%%%%%%%%%%%%%%%%%%%%%%%%%%%%%%%%%%%%%%%%%%%%%%%%%%%%%%%%%%%%%%%%%%%%%%%%%%
	%применим колонтитул к стилю страницы
\pagestyle{fancy} 
	%очистим "шапку" страницы
\fancyhead{\LaTeX} 
	%слева сверху на четных и справа на нечетных
\fancyhead[L]{\labauthors} 
	%справа сверху на четных и слева на нечетных
\fancyhead[R]{Иродов \textnumero1.136} 
	%очистим "подвал" страницы
\fancyfoot{} 
	% номер страницы в нижнем колинтуле в центре
\fancyfoot[C]{\thepage} 
%%%%%%%%%%%%%%%%%%%%%%%%%%%%%%%%%%%%%%%%%%%%%%%%%%%%%%%%%%%%%%%%%%%%%%%%%%%%%%%
\usepackage{float}
\usepackage[mode=buildnew]{standalone}
\usepackage{tikz} 
% \usepackage{subcaption}
\usepackage{tikz,csvsimple,physics}
\makeatother
%\usepackage{booktabs}
\usepackage{pgfplots, pgfplotstable}
\usepackage[outline]{contour}
\usepackage{tocloft}
\renewcommand{\cftsecleader}{\cftdotfill{\cftdotsep}}


% вот этот удивительный мир

\begin{document}
\begin{center}
    \textbf{Домашняя работа \textnumero 3}
\end{center}

\textit{Дано: $\textbf{u}$ -- скорость газа относительно ракеты, $\mu$ -- расход газа, показать, что уравнение движения $ma = F - \mu u$ }

\vspace{2em}
Рассмотрим изменение системы за малый промежуток времени $dt$.

В момент времени $t$ импульс ракеты:
\[
p(t) = mv.
\]

За время $dt$ масса ракеты изменяется на $dm$ (где $dm < 0$), скорость --- на $dv$. 

\vspace{1em}
Масса выброшенного газа: $|dm| = -dm = \mu\,dt$.

\vspace{1em}
Скорость газа относительно неподвижной системы отсчёта: $v - u$.

\vspace{1em}
Импульс в момент времени $t + dt$:
\[
p(t+dt) = (m + dm)(v + dv) + (-dm)(v - u).
\]

Раскрываем скобки и пренебрегаем членом второго порядка малости $dm\,dv$:
\[
p(t+dt) = mv + m\,dv + v\,dm + dm\,dv - v\,dm + u\,dm \approx mv + m\,dv + u\,dm.
\]

Изменение импульса:
\[
dp = p(t+dt) - p(t) = m\,dv + u\,dm.
\]

По второму закону Ньютона:
\[
dp = F\,dt.
\]

Получаем:
\[
m\,dv + u\,dm = F\,dt.
\]

Делим на $dt$:
\[
m\,\frac{dv}{dt} + u\,\frac{dm}{dt} = F.
\]

Учитываем, что $\displaystyle\frac{dv}{dt} = a$ и $\displaystyle\frac{dm}{dt} = -\mu$:
\[
ma + u(-\mu) = F.
\]

Отсюда получаем уравнение движения ракеты:
\[
\boxed{ma = F - \mu u.}
\]

\end{document}